% Generated draft: requires mathematical and source review.
\documentclass[11pt]{article}
\usepackage[margin=25mm]{geometry}
\usepackage{fontspec}
\setmainfont{Latin Modern Roman}
\usepackage{amsmath,amssymb,mathtools,unicode-math}
\setmathfont{Latin Modern Math}
\usepackage{enumitem,needspace,xurl,hyperref}
\setlength{\parindent}{0pt}
\setlength{\parskip}{5pt}
\setlength{\emergencystretch}{3em}
\newcommand{\N}{\mathbb{N}}
\newcommand{\Z}{\mathbb{Z}}
\newcommand{\Q}{\mathbb{Q}}
\newcommand{\R}{\mathbb{R}}
\newcommand{\C}{\mathbb{C}}
\newcommand{\F}{\mathbb{F}}
\newcommand{\A}{\mathbb{A}}
\newcommand{\PP}{\mathbb{P}}
\newcommand{\E}{\mathbb{E}}
\newcommand{\eps}{\varepsilon}
\newcommand{\dd}{\,\mathrm{d}}
\begin{document}
\section*{UnsolvedMath ID: 30001730}
\textbf{Source problem code:} catalogue entry
\subsection{Definitions and mathematical statement}
Let \(T\) denote the tablet called Plimpton 322 in the supplied record. Write \(R_1,\ldots,R_{15}\) for its fifteen successive rows, and refer to its columns in their successive order. These symbols refer to the historical tablet; the supplied record does not provide a transcription, numerical entries, a column count, or conventions for interpreting or correcting entries.

The source poses the following historical and mathematical reconstruction questions:
\begin{enumerate}
\item Are the successive columns of \(T\) generated by a unique algorithm?
\item Are the fifteen successive rows \(R_1,\ldots,R_{15}\) generated by a unique algorithm?
\item What tools were available to Old Babylonian scribes?
\end{enumerate}

Here ``algorithm,'' ``generated,'' ``unique,'' and ``tools'' retain their unresolved source meanings. In particular, the record specifies neither a class of admissible procedures, their inputs and permitted operations, nor an equivalence relation under which two procedures count as the same algorithm. It also does not distinguish uniqueness of a historically used procedure from uniqueness among procedures compatible with the surviving evidence. The wording does not determine whether generation concerns the numerical contents, their ordering, or both. Consequently, these are research questions requiring historical evidence and further specifications, rather than a fully specified mathematical uniqueness conjecture. The tools question likewise requires clarification of the relevant scribal setting and what counts as a tool.
\subsection{Short English statement}
The source asks whether unique algorithms generated the successive columns and the fifteen successive rows of Plimpton 322, and what tools Old Babylonian scribes could use. A precise uniqueness claim requires a tablet transcription, a specified class of procedures, a meaning of algorithmic equivalence, and criteria connecting a proposed procedure to historical practice; these are not supplied.
\subsection{Sources}
\begin{itemize}
\item Local registry: \texttt{data/unsolvedmath/problems.json}, supplied record with permanent UnsolvedMath ID \(30001730\), title ``Generation Algorithms for Plimpton 322,'' exact source problem code \texttt{catalogue entry}; collection: Oberwolfach Reports Open Problems. The record supplies the year 2011 and names no individual proposer.
\item Supplied report citation: \emph{Mini-Workshop: History of Numerical and Graphical Tables} (2011), \emph{Oberwolfach Reports}, pp.~639--689, DOI \texttt{10.4171/owr/2011/12}. Source URL and supplied literature-source link, labeled ``Oberwolfach Report 12/2011'': \url{https://doi.org/10.4171/owr/2011/12}. No specific page within that range is supplied for these questions.
\item Supplied literature assessment, checked on \texttt{2026-08-21}: the record labels the problem open and states that the questions about unique generating algorithms and scribal tools have no definitive resolution in the sources located by its curation. This is supplied dated curation, not an independently verified literature assessment; the record explicitly requests independent verification of time-sensitive status and cited claims.
\item Supplied statement-verification metadata labels the statement \texttt{corrected\_verified} and says that the official report confirms the historiographical question and that an extraction label was removed. These are assertions of the supplied registry, not independent verification performed for this draft.
\end{itemize}
\end{document}
